\documentclass[12pt,letterpaper]{article}
\usepackage[T1]{fontenc}
\usepackage{amsmath,amsthm}
\usepackage{newtxtext,newtxmath}
\usepackage[letterpaper,margin=1in]{geometry}
\usepackage{booktabs,longtable,array}
\usepackage[hidelinks]{hyperref}

\newcommand{\F}{\mathbb{F}}
\newcommand{\AGL}{\operatorname{AGL}}
\newcommand{\PAut}{\operatorname{PAut}}
\newcommand{\RS}{\operatorname{RS}}

\title{Revision log for\\
``Conjectural permutation decoding of some AG codes''\\[4pt]
\large Corrected and expanded version, Draft 2}
\author{D.~Joyner, ChatGPT (OpenAI), Claude (Anthropic)}
\date{November 12, 2026}

\begin{document}
\maketitle

\section*{Scope}
\noindent This log belongs to the corrected and expanded version of the paper.
It is kept in the folder \texttt{perm-decoding} of the repository
\begin{center}
\url{https://github.com/wdjoyner/qrcodes-hyperelliptic}.
\end{center}
The files involved are these.
\begin{center}\small
\begin{tabular}{@{}l>{\raggedright\arraybackslash}p{0.62\textwidth}@{}}
\toprule
Name used here & File\\
\midrule
2005 note & \texttt{perm\_decoding3.tex}, the source of D.~Joyner's note in
\emph{ACM SIGSAM Bulletin} 39(1) (2005), 26--32\\
Step 1 & \texttt{perm\_decoding\_step1\_mathematical\_revision.tex}, by
D.~Joyner and ChatGPT\\
Intermediate draft & \texttt{perm-decoding\_revision-2026.tex}\\
Draft 1 & \texttt{paper/perm\_decoding\_step2\_revision.tex} as committed on
November 9, 2026 (commit \texttt{188d526})\\
Draft 2 & the same file as revised on November 12, 2026\\
\bottomrule
\end{tabular}
\end{center}
The file name of the paper was kept in Draft~2 so that existing links to it stay
valid. The draft number appears on the title page.

\medskip
\noindent\textbf{How this log was made.} No revision log was kept before
Draft~2. Section~\ref{sec:draft1} is a reconstruction. Claude wrote it on November
12, 2026 by comparing the files listed above and by reading the Acknowledgments and
the \texttt{README.md} of the folder. The referee report on Step~1 and ChatGPT's
response to it are not in the repository and were not consulted. Where
Section~\ref{sec:draft1} says who found something, the source is the
Acknowledgments of the paper. Section~\ref{sec:draft2} is a direct record of the
changes made in Draft~2.

\medskip
\noindent\textbf{Numbering.} The 2005 note numbers its lemmas, propositions,
examples and conjectures with one counter (Lemma~1 to Conjecture~7), and its remark
separately. Unless stated otherwise, all other numbers in this log are Draft~2
numbers. Section~\ref{sec:conc} gives the concordance between Draft~1 and Draft~2.

%======================================================================
\section{Draft 1: from the 2005 note to the corrected and expanded version
(reconstructed)}\label{sec:draft1}

\subsection{Stages}
\begin{enumerate}
\item \emph{Step 1.} ChatGPT revised the mathematics of the 2005 note. The result
had the title ``Permutation decoding of codes from $y^2=x^p-x$: corrected results,
counterexamples, and partial decoding'' and the authors D.~Joyner and ChatGPT.
\item \emph{Referee report.} Claude refereed Step~1.
\item \emph{Intermediate draft.} ChatGPT's response to the report and the
referee's corrections were combined. Claude was added as third author.
\item \emph{Draft 1.} The original title was restored, with the subtitle
``Corrected and expanded version'', and the computations for $p=23$ were added.
\end{enumerate}

\subsection{Format and front matter}
\begin{enumerate}
\item The paper was reset in 12pt Times with one-inch margins, numbered by
section, with \verb|\label| and \verb|\ref| throughout.
\item The abstract was rewritten to state the corrected results.
\item A table of contents, a background section, a literature review (before and
after 2007), a summary, a list of open problems, a glossary and two appendices
were added. The appendices hold the GUAVA session of the 2005 note and a
verification script.
\item The bibliography grew from 30 entries (2005 note) to 32 (Step~1) to 59
(Draft~1).
\item The GUAVA session of the 2005 note was kept, with two typos corrected:
\texttt{Length(P)} became \texttt{Length(Pts)}, and the closing \texttt{if;} became
\texttt{fi;}.
\end{enumerate}

\subsection{Statements of the 2005 note that were corrected}\label{sec:errors}
Each of these corrections is flagged in the body of the paper at the place given.
\begin{enumerate}
\item \emph{Point labels.} The 2005 note takes $P_1=(1:0:1)$ and says its
stabilizer is generated by $(x,y)\mapsto(x,-y)$, $(a^2x,ay)$ and $(x+1,y)$. The
translation does not fix that point. These maps generate the stabilizer of the
point at infinity. The paper reads $P_1$ as $P_\infty$ and says so (Remark~4.1).
\item \emph{Order of the stabilizer over $\F_p$.} The 2005 note gives $2p(p-1)$.
The correct order is $p(p-1)$. The value $2p(p-1)$ is the order of the stabilizer
over $\F_{p^2}$ (Section~4.1, after equation~(4.5)).
\item \emph{Action of the hyperelliptic involution.} The proof of Lemma~2 of the
2005 note says that $\iota$ acts trivially on the monomial basis of $L(D)$. It
acts by $-1$ on the monomials that involve $y$ (proof of Proposition~4.5).
\item \emph{Lifting of M\"obius transformations.} The paper states and proves the
criterion that a class in $\mathrm{PGL}_2(p)$ lifts to an $\F_p$-automorphism
exactly when its determinant is a square (Section~4.1). According to the
Acknowledgments, the referee found an error in the lift criterion of Step~1.
\item \emph{Ample divisors and separation of points.} The 2005 note says that an
ample divisor $D$ has $L(D)$ separating points. On a curve this is false in
general (Remark~5.2).
\item \emph{Kernels.} The description of the kernel of the map from the
automorphism group to the permutation group of the code was corrected
(Proposition~5.1).
\item \emph{Permutation representations.} Proposition~4 of the 2005 note says
that, under degree hypotheses, the representation on $L(D)$ is equivalent to one
by permutation matrices. This is false. Only the weaker statement holds that
$L(D)$ is isomorphic to a submodule of a permutation module
(Proposition~5.1(3), Example~5.3).
\item \emph{Group-action convention.} The convention $g\cdot f=f\circ g^{-1}$ was
fixed and used throughout (equation~(3.1)).
\end{enumerate}

\subsection{The two conjectures}
\begin{enumerate}
\item Conjecture~6 of the 2005 note, called (C1) in the paper, is false as a
universal statement. At $p=37$ and $k=17$ no information set works
(Theorem~7.8). The counterexample is already in Step~1.
\item Conjecture~7 of the 2005 note, called (C2), is false as a universal
statement. Step~1 gave one counterexample, at $p=7$ (its Theorem on the
$[84,42,\ge40]$ code). Draft~1 replaced it by a general obstruction for regular
groups (Theorem~8.1) and an infinite family (Corollary~8.2). The case $p=7$
became Example~8.3.
\end{enumerate}

\subsection{Changes between Step 1 and Draft 1}
\begin{enumerate}
\item The original title was restored, and Claude was added as third author.
\item The Background and Literature review sections, the Glossary and
Appendix~B were added.
\item Proposition~7.5 (exhaustive computations for $p\le23$), Remark~7.6 and
Example~7.7 (failure at $p=17$ despite the counting bound) were added. The
intermediate draft covered $p\le19$, and Draft~1 extended this to $p=23$.
\item Theorem~8.1, Corollary~8.2 and the table after it were added, as described
above. According to the Acknowledgments, Claude proposed the logarithmic
obstruction, and ChatGPT found the choice $m=n/2+g-1$ and the information set of
Example~7.7.
\item Problems~10.2 and~10.4 and the references to the recent literature on
permutation decoding of AG codes were added.
\item The corrections 1, 4, 6 and 8 of Section~\ref{sec:errors} were made at this
stage, according to the Acknowledgments.
\end{enumerate}

%======================================================================
\section{Draft 2: revision for the style guide (November 12, 2026)}
\label{sec:draft2}

\subsection{Scope}
Claude (Claude Fable~5.2) brought the paper into line with D.~Joyner's style
guide, following his instructions of November 12, 2026. The guide is the file
\texttt{joyner-math-paper-style-guide.md} at the top level of the repository
(Draft~2 of November 11, 2026). Those instructions forbade
mathematical changes. No statement, proof, hypothesis or status label was altered.
All new mathematical text was confined to what the guide requires, and it is
listed in Section~\ref{sec:newtext} so that it can be refereed. Semicolons in
existing text were left in place. Rule~3.8 (no semicolons) was applied to new
text only. The guide has no paper-specific rule (its Section~9) for this paper.

Before the changes, the paper compiled with four errors, all at the title of
Example~7.4. The optional argument of that example contained unprotected square
brackets, so the title was cut short and the first sentence of the example was set
in math mode. There were also 16 hyperref warnings, caused by math in five section
titles, one overfull box (the display~(4.4), by 53pt) and two underfull boxes in
the bibliography. After the changes, the paper compiles with no errors, no
warnings, no undefined references and no overfull boxes. One underfull box remains,
in the bibliography entry \texttt{Go03}.

\subsection{Defaults used for the rules marked [DECIDE]}
\begin{enumerate}
\item Document class: the existing class, \texttt{article}, was kept.
\item Times package: \texttt{newtxtext} and \texttt{newtxmath}.
\item Glossary placement: notation near the beginning (Subsection~3.7), and the
glossary at the end, just before the Acknowledgments.
\item AI systems on the author line: the instructions did not mention the author
line, so it was left as it was (D.~Joyner, ChatGPT, Claude). The Acknowledgments
also credit both systems and point to this log.
\end{enumerate}

\subsection{Changes}
\paragraph{Typesetting.}
\begin{enumerate}
\item The class and \texttt{geometry} options now name letter paper. The
conditional fallback to \texttt{mathptmx} was removed, and \texttt{newtxtext} and
\texttt{newtxmath} are loaded directly.
\item The title of Example~7.4 is now enclosed in braces inside the optional
argument.
\item The display~(4.4) was too wide by 53pt. It is now set on two lines. Its
content and its number are unchanged.
\item The five section and subsection titles that contain math now use
\verb|\texorpdfstring|, which removes the hyperref warnings.
\item In Subsection~2.2 the longer key \texttt{CLLS25} made one line too wide. The
citation was moved four words later in the same sentence, which now reads
``\dots\ and Santos gave an elementary proof [CLLS25] that \dots''.
\end{enumerate}

\paragraph{Title page, abstract, Acknowledgments.}
\begin{enumerate}
\setcounter{enumi}{5}
\item The subtitle became ``Corrected and expanded version. Draft 2'' (rule~6.2).
The title is unchanged.
\item The date line became ``Original draft: 2005. Corrected and expanded draft:
2026.'' (rule~6.2).
\item A \verb|\thanks| note on the title identifies the 2005 note and gives the
date of Draft~2 (rule~6.3).
\item The Naval Academy address was removed from the first author's footnote,
which now reads ``Original author. \texttt{wdjoyner@gmail.com}'' (rule~5.1).
\item The first sentence of the abstract, ``This AI-aided paper revisits \dots'',
was replaced by a disclosure sentence that names both AI systems and refers to
the Acknowledgments, followed by ``The paper revisits \dots'' (rule~5.2). The
rest of the abstract is unchanged.
\item The Acknowledgments gained two paragraphs. The first describes Claude's
role in Draft~2 and cites the check script by commit (rules~5.3 and~8.4). The
second says that AI review is not an independent certification and points to
this log (rules~5.4 and~5.5).
\end{enumerate}

\paragraph{Citation keys.}
\begin{enumerate}
\setcounter{enumi}{11}
\item Sixteen keys and their labels were changed to the pattern of rule~1.7, and
every \verb|\cite| was updated.
\begin{center}\small\ttfamily
\begin{tabular}{@{}ll@{\qquad}ll@{\qquad}ll@{\qquad}ll@{}}
\toprule
\rmfamily Old & \rmfamily New & \rmfamily Old & \rmfamily New & \rmfamily Old & \rmfamily New & \rmfamily Old & \rmfamily New\\
\midrule
CLLS & CLLS25 & HvLP & HvLP98 & JT & JT02 & Sti & St93\\
G & Go03 & HP & HP03 & KW & KW90 & TV & TV91\\
Gu & Gu04 & JK2 & JK03 & L & Le82 & W & We98\\
H & Ha77 & JKT & JKT08 & S & Sc90 & repo & Jo26\\
\bottomrule
\end{tabular}
\end{center}
The label of \texttt{Du87} lost its umlaut. The years in \texttt{Go03},
\texttt{JK03}, \texttt{JT02} and \texttt{JKT08} are taken from the arXiv numbers
in the existing entries. No bibliographic data was changed.
\item The entry \texttt{Jo26} (the repository) was moved next to \texttt{Jo05}.
The widest label of the bibliography is now \texttt{JLJHH89}.
\end{enumerate}

\paragraph{Section order.}
\begin{enumerate}
\setcounter{enumi}{13}
\item The Literature review now comes before the Background (rule~2). They were
Sections~3 and~2 and are now Sections~2 and~3.
\item The former Section~10 (Open problems) and Section~11 (Summary) were merged
into one Section~10, ``Summary of results and open problems'', with the summary
first (rule~2.7). The label \texttt{sec:problems} now points to the merged
section. The problems keep the numbers 10.1 to 10.4.
\item The Acknowledgments were moved from before the appendices to after the
Glossary, just before the references (rule~2.9).
\item The ``Organization'' paragraph of the Introduction was updated to match
the new order.
\end{enumerate}

\subsection{New text to be refereed}\label{sec:newtext}
Rules~2.4, 2.5, 3.2, 3.3 and~3.4 required the following additions. Every number
in N1 to N10 is checked, by exhaustive and exact computation, by the script
\texttt{check\_examples\_draft2.py} in the directory \texttt{scripts}
(commit \texttt{9e1e1ea}).
\begin{description}
\item[N1.] \textbf{Example~1.1} (\texttt{ex:intro-p7}). The $[7,3,5]$ code as an
instance of the setup. It restates facts of Example~7.4.
\item[N2.] \textbf{Example~3.1} (\texttt{ex:bg-codes}). Information set, check
positions and the parity-check matrix for the generator matrix of
Example~7.4. The parity-check matrix is new.
\item[N3.] \textbf{Example~3.2} (\texttt{ex:rs7}). Two codewords of
$\RS(\F_7,3)$ and one affine map.
\item[N4.] \textbf{Example~3.3} (\texttt{ex:bg-p1}). $\RS(\F_7,3)$ as an AG code
on the projective line.
\item[N5.] \textbf{Example~3.4} (\texttt{ex:bg-hyper}). Weierstrass points and
point counts of $y^2=x^7-x$ over $\F_7$ and $\F_{49}$.
\item[N6.] \textbf{Example~3.5} (\texttt{ex:bg-groups}). The action of
$\AGL(1,7)$ on $\F_7$. It includes a new small claim: the sum-zero submodule of
$\F_7^{\,7}$ is not a permutation module for $\AGL(1,7)$. The argument is the one
used in Example~5.3.
\item[N7.] \textbf{Example~3.6} (\texttt{ex:bg-pgl}). Orders of
$\mathrm{PGL}_2(7)$, $\mathrm{PSL}_2(7)$ and the Borel subgroup, and two maps
that are and are not in $\mathrm{PSL}_2(7)$.
\item[N8.] \textbf{Example~4.8} (\texttt{ex:curve-p7}). The case $p=7$ of
Section~4: a basis of $L(7P_\infty)$, the orders of $G_7$, $H_7$, $G_{49}$ and
$H$, the orbit sizes, and $\ell(D_2)=82$. It also serves as the example for
Definition~4.3 (large group).
\item[N9.] \textbf{Example~6.5} (\texttt{ex:pd-translations}). It serves as the
example for Definition~6.2. It contains two new small claims. The seven
translations form a $2$-PD-set for $\RS(\F_7,3)$ relative to
$\mathcal I=\{0,1,2\}$. No set of three permutations is a $2$-PD-set for that
code.
\item[N10.] \textbf{Example~9.1} (\texttt{ex:cost}). Operation counts for
$p=7$ and for $\RS(\F_{23},11)$.
\item[N11.] \textbf{Table of notation} (Subsection~3.7), with an orienting
paragraph (rule~2.5). Each row points to the place where the symbol is
introduced.
\item[N12.] \textbf{Orienting paragraphs} (rule~3.2) at the start of
Sections~1, 2, 4, 5, 6, 7, 8, 9 and~10, of Subsections~2.1, 2.2, 3.1 to 3.6, 4.1,
4.2, 5.1, 7.1, 7.2, 10.1 and~10.2, and of the Glossary. Two sentences were added
to the existing opening paragraph of Section~3. The paragraph of
Subsection~10.1 makes a statement about status: the permutation group in item~1
of the summary is cited from the literature, item~3 rests in part on the computer
search of Proposition~7.5, and everything else is proved in the paper. A referee
should confirm that this describes the existing results correctly.
\item[N13.] Seven new labels for the subsections of Section~3
(\texttt{subsec:bg-codes} and so on, and \texttt{subsec:notation}).
\end{description}

\subsection{Concordance of numbers}\label{sec:conc}
Every \verb|\label| of Draft~1 was kept. The table was generated by comparing
the \texttt{.aux} files of the two drafts. All new examples were placed either
in sections that had no numbered statements or at the end of a section, so no
theorem, proposition, lemma, corollary, remark, example or problem of Draft~1
changed its number.
\begin{center}\small
\begin{tabular}{@{}lll@{}}
\toprule
Draft 1 & Draft 2 & Label or title\\
\midrule
Section 2 & Section 3 & \texttt{sec:background}\\
Subsections 2.1--2.6 & Subsections 3.1--3.6 & (Background)\\
Equation (2.1) & Equation (3.1) & \texttt{eq:convention}\\
Section 3 & Section 2 & \texttt{sec:lit}\\
Subsections 3.1, 3.2 & Subsections 2.1, 2.2 & (Literature review)\\
Section 10 & Subsection 10.2 & Open problems (\texttt{sec:problems} is now Section 10)\\
Section 11 & Subsection 10.1 & Summary\\
\bottomrule
\end{tabular}
\end{center}
New numbers in Draft~2: Examples~1.1, 3.1--3.6, 4.8, 6.5 and~9.1, and
Subsection~3.7.

\subsection{Rules not applied, and why}
D.~Joyner approved items 1 to 5 before the changes were made.
\begin{enumerate}
\item Rule~6.4 for the corrections of the 2005 note in the body (Remark~4.1, the
note on the stabilizer order in Section~4.1, the proof of Proposition~4.5,
Remark~5.2, Subsection~5.1, the paragraph ``Relation to the original note'').
They are mathematical statements about a published paper, and removing them
would change the mathematics.
\item Rule~3.3 in Section~2 (Literature review), Section~10 (Summary of results
and open problems), the appendices and the Glossary. An example there would add
nothing.
\item Rule~1.7 for the key \texttt{GUAVA}, which is software cited by name.
\item Rule~8.3 as a change of file name. The file is still
\texttt{perm\_decoding\_step2\_revision.tex}.
\item Rule~1.6, second sentence (unnumbered subsections). Rule~3.8 in existing
text. Rules~3.6 and~3.7, where no violation was looked for. Rule~6.6, since no
status label was added.
\item Rule~2.3 asks that the literature review be split at the date of the
original paper, which is 2005. The existing split is at 2007. Moving it would mean
rewriting the paragraphs on work of 2006 and 2007, so it was left for
D.~Joyner to decide.
\item Rule~1.5 for citations typed by hand. None was found.
\item Rules~8.5 and~8.6 concern the repository. It has no license file. The
copyright position of this corrected version relative to the published 2005 note
is not recorded. Neither point was acted on.
\end{enumerate}

\subsection{Not verified}
\begin{enumerate}
\item No cited source was consulted in Draft~2 (rule~4.3). See
Section~\ref{sec:notconsulted}.
\item The computations in the body were rerun in Draft~2 (rule~4.4) and
confirmed. The program \texttt{pd\_affine.py} was rerun for every prime
$5\le p\le23$. Its output agrees with the stored files in \texttt{results}
apart from the timings, and the table of Proposition~7.5 agrees with that output
entry by entry. The order $42$ of the permutation group of the $[7,3,5]$ code was
confirmed by a search of all of $S_7$ in the check script. The other numerical
statements of the body (Example~7.7, Theorem~7.8, Remark~7.9, the table after
Corollary~8.2 and Example~8.3) were rechecked with the new script
\texttt{check\_body\_numbers.py}, and all were confirmed. No hand proof in
the body was rechecked.
\item The new text N1 to N13 was checked by its author and by the check script,
not by an independent referee.
\item Section~\ref{sec:draft1} is a reconstruction, as explained under Scope.
\item The versions of ChatGPT and Claude that did the work before Draft~2 are not
recorded in the repository.
\end{enumerate}

%======================================================================
\section{AI usage}\label{sec:ai}
\begin{enumerate}
\item \emph{ChatGPT (OpenAI): first reviser.} It prepared Step~1 and responded to
the referee report. The Acknowledgments of the paper list its contributions.
\item \emph{Claude (Anthropic): referee, co-author and reviser.} It refereed
Step~1, carried out the exhaustive computations of Proposition~7.5, wrote the
scripts in \texttt{scripts/}, and prepared Draft~1, including the Background,
the Literature review and the Glossary. The Acknowledgments list its
contributions.
\item \emph{Claude (Claude Fable~5.2): style reviser, Draft~2.} It made the
changes in Section~\ref{sec:draft2}, wrote the new text N1 to N13 and the check script
for it, and wrote this log.
\end{enumerate}
D.~Joyner directed the work and is responsible for the final content. AI review is
not an independent certification of the results.

\section{Sources not consulted}\label{sec:notconsulted}
Rule~4.3 asks that every citation be checked against its source. In Draft~2 no
cited source was consulted. The Acknowledgments of Draft~1 say that Claude
verified the bibliography at that stage, but the repository holds no record of
which sources were read. In particular, the following precise citations were not
checked in Draft~2.
\begin{enumerate}
\item From Stichtenoth~\cite{St93}: Theorems~I.5.17, II.2.2 and~V.2.3,
Proposition~VI.4.1 and \S VII.3.
\item From Hartshorne~\cite{Ha77}: Corollary~IV.3.2.
\item From MacWilliams and Sloane~\cite{MS77}: Chapters~11 and~16.
\item The theorem of Kn\"orr and Willems~\cite{KW90} used in Theorem~7.1.
\end{enumerate}

\section{Original and new results}\label{sec:map}
Rule~6.5 asks that original results be separated from results first established
in the revision. Draft~2 adds no new results apart from the small claims in N6
and N9.
\begin{longtable}{@{}p{0.40\textwidth}p{0.55\textwidth}@{}}
\toprule
2005 note & Draft~2\\
\midrule\endhead
Definition of a large group & Definition 4.3\\
Lemma 1 (large groups ramify) & Lemma 4.4\\
Monomial basis of $L(mP)$, quoted from Stichtenoth & Lemma 4.2 (with proof)\\
Lemma 2 (semisimplification) & Proposition 4.5 (statement on $\iota$ in the proof corrected)\\
Proposition 3 (orbits over $\F_{p^2}$, proof omitted) & Proposition 4.6 (with proof)\\
Kernel and separation of points (Section 3.1) & Proposition 5.1, Remark 5.2 (corrected)\\
Proposition 4 (permutation matrices) & false, replaced by Proposition 5.1(3) and Example 5.3\\
Syndrome lemma, stated without proof & Lemma 6.1 (with proof)\\
Example 5 (the $[7,3,5]$ code) & Example 7.4\\
Conjecture 6 & (C1), refuted by Theorem 7.8\\
Conjecture 7 & (C2), refuted by Corollary 8.2\\
--- & New in the revision: Theorem 6.3 (counting bound, after Gordon),
Theorem 7.1 (Reed--Solomon identification), Theorem 7.3 (two errors),
Proposition 7.5, Remark 7.6 and Example 7.7 (computations for $p\le23$),
Theorem 7.8, Theorem 8.1, Corollary 8.2, Example 8.3, Example 5.3,
Problems 10.1--10.4\\
\bottomrule
\end{longtable}

\begin{thebibliography}{KW90}
\bibitem[Ha77]{Ha77} R.~Hartshorne, \emph{Algebraic Geometry}, Springer, 1977.
\bibitem[KW90]{KW90} R.~Kn\"orr and W.~Willems, The automorphism groups of
generalized Reed--Muller codes, \emph{Ast\'erisque} 181--182 (1990), 195--207.
\bibitem[MS77]{MS77} F.~J.~MacWilliams and N.~J.~A.~Sloane, \emph{The Theory of
Error-Correcting Codes}, North-Holland, 1977.
\bibitem[St93]{St93} H.~Stichtenoth, \emph{Algebraic Function Fields and Codes},
Springer, 1993 (2nd ed.\ 2009).
\end{thebibliography}
\end{document}
